\abstract

\keywords{интеллектуальный анализ данных, реляционные базы данных, Text-to-SQL, большие языковые модели, RAG, SQL, микросервисная архитектура, самокоррекция, векторное хранилище, ChromaDB}

Объектом разработки является интеллектуальный модуль интерактивного анализа данных в реляционных базах данных. Цель работы состоит в разработке программного решения, позволяющего пользователям без специальной подготовки формулировать запросы к данным на естественном языке и получать результат без ручного написания SQL.

В работе рассмотрены современные подходы к построению интерфейсов естественного языка к базам данных, методы Text-to-SQL и применение больших языковых моделей совместно с RAG-архитектурой. Разработанное решение использует микросервисную архитектуру, векторное хранилище для поиска релевантного контекста схемы данных и агентный механизм самокоррекции SQL-запросов на основе обратной связи от СУБД.

Практическая значимость работы заключается в снижении порога доступа к корпоративным данным, уменьшении нагрузки на технических специалистов и повышении автономности аналитических сценариев. Результаты могут применяться в организациях, использующих реляционные базы данных и нуждающихся в интерактивном доступе к информации.
